\chapter{Introduction}
\label{chap:introduction}

This chapter motivates repository maintenance as a dependency-security concern,
states the research question, and defines the scope and contributions of the
thesis.

\section{Motivation and Problem Statement}
\label{sec:intro-motivation}

Modern software systems combine first-party code with third-party and
open-source dependencies. Studies across Java, Python, Ruby, and npm describe
applications that depend substantially on externally maintained packages,
including transitive dependencies
\cite{prana2021outsight,miller2025abandonment}. Here, open-source software means
software distributed under a license that complies with the Open Source
Definition \cite{osiOpenSourceDefinition}; Section~\ref{sec:background-supply-chain-security}
gives the licensing criteria.

Depending on external components transfers part of the maintenance burden to
people outside the downstream project. If maintainers no longer address a
disclosed vulnerability, the downstream project must fork the component, vendor
a private patch, or migrate instead of applying a normal version update. Until
then, the dependency remains exposed without an upstream fix.
Secure-development guidance therefore treats
dependency upkeep as a supply-chain control. The NIST Secure Software
Development Framework includes acquiring and maintaining well-secured
third-party and open-source components \cite{souppaya2022ssdf}. SLSA addresses
artifact and build-process integrity and provenance
\cite{slsa2025spec}.\kiref{Claude Opus 5/1}

Miller et al. found that \(15\%\) of widely used npm packages became abandoned
within six years. Only \(18\%\) of the exposed downstream projects removed the
abandoned dependency \cite{miller2025abandonment}. Their earlier interview study
found that affected developers often lack clear response strategies
\cite{miller2023wingingit}. Delayed maintenance also prolongs security exposure.
Across dependencies in 450 Java, Python, and Ruby projects, publicly known
vulnerabilities persisted throughout a one-year observation period at mean rates
of \(78.4\%\), \(97.7\%\), and \(66.4\%\), respectively. Resolved
vulnerabilities took three to five months \cite{prana2021outsight}.\kiref{Claude Opus 5/10}

Several incidents illustrate the consequences. The 2016 \texttt{left-pad}
unpublishing disrupted many direct and transitive dependants
\cite{npm2016leftpad}. The 2017 Equifax breach exploited an Apache Struts
vulnerability whose available patch had not been applied, and Equifax later
agreed to a settlement of up to
\$700 million \cite{equifax2017incident,ftc2019equifax}. The XZ Utils backdoor
made maintainer trust an attack surface through a social-engineering takeover
\cite{openssf2024xzalert}.

An inactive repository is not necessarily insecure or abandoned; a project may
be stable, paused, or maintained through channels not visible in public data.
The examples nevertheless show how dependency maintenance can affect downstream
security exposure. This thesis uses repository inactivity as an operational
proxy for elevated maintenance risk and examines whether public signals can
estimate that risk over a 12-month horizon.

\section{Research Question and Objectives}
\label{sec:intro-research-question}

The central research question of this thesis is:

\begin{quote}
To what extent can a machine-learning-based scoring tool, using only publicly observable repository health and maintenance indicators, predict 12-month OSS dependency inactivity, and how can this inactivity-risk score be combined with parallel security-practice signals to support risk-based dependency triage in secure dependency management?
\end{quote}

To answer this research question, the thesis pursues the following objectives:

\begin{enumerate}
    \item Define an operational inactivity label for open-source repositories within a 12-month prediction horizon.
    \item Construct a reproducible dataset from publicly observable repository metadata and dependency-related sources.
    \item Engineer observation-time activity and maintenance features for the inactivity model, and surface selected security-practice signals as parallel triage context rather than as inputs to the inactivity prediction.
    \item Train and compare simple baselines and machine-learning models for inactivity-risk prediction.
    \item Evaluate the models using discrimination and calibration metrics that are suitable for risk scoring.
    \item Demonstrate how the trained model can be applied per repository to a set of dependency repositories to support risk-based triage.
\end{enumerate}

The research follows a Design Science Research-oriented approach and develops an
inactivity-risk scoring workflow for open-source dependencies
\cite{hevner2004design,peffers2007design}. The empirical evaluation measures
discrimination and calibration. The \emph{demonstrator} is a worked application
of the artifact to a real case: Chapter~\ref{chap:demonstrator} applies the
trained model to the dependency repositories of an existing open-source project.
It shows integration into dependency monitoring and behavior outside the sampled
evaluation set, while
Chapter~\ref{chap:evaluation} provides the measurement of predictive
accuracy. This distinction defines the role of the \emph{demonstrator}
throughout the thesis.\kiref{Claude Opus 5/1}

\section{Scope and Assumptions}
\label{sec:intro-scope}

The scope covers dependencies with publicly observable source repositories and
package metadata that can be associated with an open-source license through
GitHub's SPDX-based license information \cite{githubRestLicenses}. It excludes
private repositories, closed-source vendor components, and packages that cannot
be mapped to a source repository (Section~\ref{sec:method-data-sources}).

The thesis covers only projects developed on \emph{GitHub}. Both training data
and observation-time features are reconstructed from its public event stream
(Section~\ref{sec:method-data-sources}). Projects on GitLab, Codeberg,
self-hosted forges, or mailing-list workflows such as that of the Linux kernel
are outside the scope. This non-random exclusion limits generalization to
open-source software as a whole (Section~\ref{sec:validity-external}). A GitHub
mirror can also appear inactive although development continues
elsewhere.\kiref{Claude Opus 5/1}

The prediction target is repository inactivity within
\( [t, t + 12\ \text{months}] \). The label uses four activity checks: human
commit activity, active contributors, releases or package publications, and
merged pull requests. At least two must hold for the repository to count as
maintained (Section~\ref{sec:method-definitions}), building on revision activity
as a survival signal \cite{robinson2022survival}. Merged pull requests serve as a
proxy for maintainer interaction. Issue activity and responsiveness are
predictors only and require caution because bots can distort response measures
\cite{hasan2023firstresponse}.

Public signals provide incomplete evidence about future inactivity, and health
indicators can have different meanings across projects
\cite{yehudi2023contextfree}. The model output therefore supports triage but does
not automate the decision.

\section{Contributions}
\label{sec:intro-contributions}

The thesis makes three contributions:

\begin{enumerate}
    \item an operational 12-month inactivity-risk label and a reproducible workflow for constructing it;
    \item an empirical evaluation of inactivity prediction from public repository health and maintenance indicators, benchmarked against simple baselines on discrimination and probability calibration, with security-practice signals kept outside the inactivity model as complementary triage context;
    \item a per-repository demonstrator that turns the score into a ranked triage view.
\end{enumerate}

The score provides a reproducible way to prioritize dependencies for expert
review.

\section{Structure of the Thesis}
\label{sec:intro-structure}

Chapters~\ref{chap:background} and~\ref{chap:related-work} provide the conceptual
and empirical background. Chapters~\ref{chap:methodology}
and~\ref{chap:implementation} describe the method and its implementation.
Chapter~\ref{chap:evaluation} evaluates the models, and
Chapter~\ref{chap:demonstrator} applies them in a ranked triage view.
Chapters~\ref{chap:validity} and~\ref{chap:conclusion} discuss validity and future
work.
